%%%%%%%%%%%%%%%%%%%%%%%%%%%%%%%%%%%%%%%%%%%%%%%%%%%%%%%%%%%%%%%%%%%%%%%%%%%%%%%%
%%  EINSTELLUNGEN – die einzige Datei, in der du deine Angaben einträgst.
%%  Wird von ALLEN Hauptdateien (Endabgabe + Zwischenabgaben) eingelesen.
%%  Entspricht dem Block "Die Angaben" der offiziellen Gerüste
%%  pp-portfolio-geruest-2026-2027.tex / pp-portfolio-geruest-einlage-2026-2027.tex
%%%%%%%%%%%%%%%%%%%%%%%%%%%%%%%%%%%%%%%%%%%%%%%%%%%%%%%%%%%%%%%%%%%%%%%%%%%%%%%%

%%%%%%%%%%%%%%%%%%%%%%%%%%%%%%%%%%%%% Die Angaben %%%%%%%%%%%%%%%%%%%%%%%%%%%%%%%
\newcommand{\Semester}{WH 2026-2027}
\newcommand{\Projekt}{Projekt, Titel des Portfolios}          % z.B. Titel eures Gruppenprojekts
\newcommand{\Schwerpunkt}{Ihr Schwerpunktthema}               % DEIN Schwerpunkt (≠ Partner/in!)
\newcommand{\IhrName}{Moritz Kautzsch}                         % Vorname Nachname
\newcommand{\Matrikelnummer}{Ihre Matrikelnummer}
\newcommand{\Studiengang}{Ihr Studiengang (Name und Stufe)}   % z.B. Physik B.Sc., 5. Semester
\newcommand{\emailAdresse}{Ihre e-Mail-Adresse}               % Uni-Adresse empfohlen
\newcommand{\dasDatum}{Das Datum der Abgabe}                  % z.B. 15.\,03.\,2027

%% Weitere Angaben, die in mehreren Einlagen gebraucht werden
\newcommand{\Partner}{Name Partnerin/Partner}                  % für Abgrenzung der Schwerpunkte
\newcommand{\PartnerSchwerpunkt}{Schwerpunkt Partnerin/Partner}
\newcommand{\Laborbuchlink}{https://link-zum-elektronischen-laborjournal}  % Pflicht in Einlage 3!
\newcommand{\Protokolldatum}{TT.MM.2026}   % Einlage 3: gewählter Tag (≠ Tag der Partnerin/des Partners)
\newcommand{\PaperKey}{VORLAGE-Artikel}        % Einlage 2: bib-Key deines Papers in literatur.bib
\newcommand{\Vortragsdatum}{28.\,01.\,2027}    % Einlage 5
\newcommand{\FremdVortragKey}{VORLAGE-Vortrag} % Einlage 6: bib-Key des anderen Vortrags

%%%%%%%%%%%%%%%%%%%%%%%%%%%%%%%%%%%%%%%%%%%%%%%%%%%%%%%%%%%%%%%%%%%%%%%%%%%%%%%%
%%  MODUS-SCHALTER
%%
%%  \Abgabefalse  = Entwurf: grüne Hilfen sichtbar, Kontrollseite mit
%%                  Umfangsmessung, Platzhalter rot markiert.
%%  \Abgabetrue   = Abgabe: grüne Farbe -> schwarz (wie offizielle Zeile 15/16),
%%                  alle Hilfen + Kontrollseite + Umfangsboxen verschwinden.
%%                  Übrig gebliebene \platzhalter{} bleiben ROT sichtbar und
%%                  erzeugen eine Warnung -> dann ist noch etwas offen!
%%
%%  \Hinweisefalse blendet im Entwurf nur die grünen Hilfstexte aus
%%                  (nützlich, um den Umfang realistisch zu messen).
%%%%%%%%%%%%%%%%%%%%%%%%%%%%%%%%%%%%%%%%%%%%%%%%%%%%%%%%%%%%%%%%%%%%%%%%%%%%%%%%
\newif\ifAbgabe
\newif\ifHinweise

\Abgabefalse      % <-- für die Abgabe auf \Abgabetrue setzen
\Hinweisetrue     % <-- Hilfstexte im Entwurf ein-/ausblenden
